\documentclass[11pt,oneside]{book}

\usepackage{fontspec}
\setmainfont{TeX Gyre Pagella}

\usepackage{unicode-math}
\setmathfont{Latin Modern Math}

\usepackage[margin=1.1in]{geometry}

\usepackage{amsmath}
\usepackage{amssymb}
\usepackage{amsthm}
\usepackage{mathtools}
\usepackage{mathrsfs}
\usepackage{stmaryrd}
\usepackage{bm}

\usepackage{microtype}

\usepackage{xcolor}

\usepackage[
colorlinks=true,
linkcolor=blue!50!black,
citecolor=blue!50!black,
urlcolor=blue!50!black
]{hyperref}

\usepackage[nameinlink]{cleveref}

\usepackage{enumitem}

\usepackage{booktabs}
\usepackage{array}
\usepackage{longtable}

\usepackage{tikz}
\usetikzlibrary{
arrows.meta,
calc,
matrix,
positioning,
decorations.pathmorphing,
fit
}

\usepackage{graphicx}

\usepackage{imakeidx}
\makeindex

\numberwithin{equation}{chapter}

\theoremstyle{definition}
\newtheorem{definition}{Definition}[chapter]
\newtheorem{example}[definition]{Example}
\newtheorem{remark}[definition]{Remark}

\theoremstyle{plain}
\newtheorem{theorem}[definition]{Theorem}
\newtheorem{lemma}[definition]{Lemma}
\newtheorem{corollary}[definition]{Corollary}
\newtheorem{proposition}[definition]{Proposition}

\DeclareMathOperator{\id}{id}
\DeclareMathOperator{\Hom}{Hom}
\DeclareMathOperator{\Obj}{Obj}
\DeclareMathOperator{\End}{End}
\DeclareMathOperator{\Aut}{Aut}
\DeclareMathOperator{\Int}{Int}
\DeclareMathOperator{\cl}{cl}

\newcommand{\Adm}{\mathsf{Adm}}
\newcommand{\State}{\mathcal{S}}
\newcommand{\Env}{\mathcal{E}}
\newcommand{\Theory}{\mathfrak{T}}
\newcommand{\Repair}{\mathcal{R}}
\newcommand{\Continue}{\mathcal{C}}
\newcommand{\Invariant}{\mathcal{J}}

\begin{document}

\begin{titlepage}

\centering

\vspace*{2.5cm}

{\Huge\bfseries
The Mathematics of Autonomous Systems\par}

\vspace{0.8cm}

{\Large
Admissibility, Repair, Continuation,\\
Representation, and Scientific Discovery\par}

\vfill

{\Large
Flyxion\par}

\vspace{0.3cm}

{\large
Independent Researcher\par}

\vfill

{\large
\today\par}

\end{titlepage}

\frontmatter

\chapter*{Dedication}

\begin{flushright}
\emph{
To everyone who believes\\
that understanding\\
is not accumulated,\\
but recursively repaired.
}
\end{flushright}

\cleardoublepage


\begin{abstract}

This monograph develops a unified mathematical framework for autonomous
systems, admissibility, recursive repair, continuation, observation,
representation, and scientific discovery.

Rather than treating cognition, computation, biology, scientific inquiry,
and distributed intelligence as independent subjects, each is formulated as
an instance of constrained evolution upon admissible state spaces.

Beginning with the notions of admissible interfaces and invariant
preservation, we construct a hierarchy of mathematical structures including
repair operators, continuation semigroups, distinction geometry,
reachability manifolds, sheaf-theoretic reconstruction, categorical
observation, and recursive theory evolution.

Within this framework, identity emerges from admissible continuation rather
than static substance, information arises from distinguishability rather
than probability alone, and scientific progress becomes recursive navigation
through the geometry of theory space.

The resulting formalism replaces numerous informal notions—including
boundaries, sanctuaries, equilibrium, optimization, and discovery—with
general mathematical objects applicable across physical, biological,
computational, and epistemological domains.

The central thesis is that autonomy is neither isolation nor independence,
but invariant-preserving interaction through admissible interfaces.
Likewise, scientific knowledge is neither a collection of isolated facts nor
an accumulation of arbitrary hypotheses, but the recursive repair of
representational structures constrained by observation and governed by
continuation.

\end{abstract}

\tableofcontents

\mainmatter

\chapter{Autonomous Cognitive Domains and Admissible Interfaces}

\section{Motivation}

Many descriptions of creative, intellectual, or contemplative work invoke
metaphors such as sanctuaries, fortresses, halls, monasteries, laboratories,
or isolated kingdoms. Although such imagery is psychologically compelling,
it obscures a deeper mathematical structure.

The purpose of this chapter is to replace metaphor with formalism.

We demonstrate that what is commonly perceived as an isolated sanctuary is
better understood as an \emph{autonomous cognitive domain}: an open dynamical
system whose interfaces satisfy admissibility constraints. Such systems are
not closed. They exchange information, resources, and computational work with
their environment while selectively rejecting perturbations that increase
structural entropy.

The resulting framework applies equally to biological organisms, software
systems, scientific collaborations, research laboratories, distributed
computing networks, and individual cognitive agents.

Consequently, architectural metaphors become manifestations of a more general
mathematical object rather than primitive explanatory concepts.

\section{State Spaces}

Let

\[
\mathcal{S}
\]

denote the internal state space of an autonomous system.

Let

\[
\mathcal{E}
\]

denote its environment.

Neither space is assumed finite, compact, linear, or even metrizable unless
explicitly stated.

The evolution of the internal configuration is represented by

\[
\Phi_t:\mathcal{S}\rightarrow\mathcal{S},
\]

while environmental evolution is represented by

\[
\Psi_t:\mathcal{E}\rightarrow\mathcal{E}.
\]

The interaction between the two is not direct.

Instead it is mediated by a boundary operator

\[
\mathcal{B}:
\partial\mathcal{S}
\rightarrow
\partial\mathcal{E},
\]

whose admissibility properties determine which external influences are
permitted to alter internal state.

Thus autonomy is not equivalent to isolation.

Autonomy is the existence of a mathematically constrained interface.

\section{The Error of Closed-System Thinking}

A common misconception is that sustained intellectual work requires complete
disconnection from the external world.

Such an assumption is impossible for any computationally active system.

Every scientific institution, every software project, every biological
organism, and every civilization continuously exchanges matter, energy,
information, or computation with its surroundings.

Therefore complete closure implies eventual computational death.

The relevant distinction is therefore not

\[
\text{closed}
\qquad
\text{versus}
\qquad
\text{open},
\]

but instead

\[
\text{indiscriminately coupled}
\qquad
\text{versus}
\qquad
\text{admissibly coupled}.
\]

This distinction forms the foundation of the present theory.

\section{Admissible Interfaces}

The defining property of an autonomous system is not impermeability but
controlled interaction.

Every computational system capable of sustained existence exchanges
information with its surroundings. The mathematical problem is therefore not
how to eliminate interaction, but how to constrain it.

\begin{definition}[Interface]
Let $\mathcal{S}$ denote an autonomous state space and let
$\mathcal{E}$ denote its environment.

An \emph{interface} is a morphism

\[
I:\mathcal{E}\longrightarrow\mathcal{S}
\]

together with a corresponding outward morphism

\[
O:\mathcal{S}\longrightarrow\mathcal{E},
\]

describing the admissible exchange of information, computation,
resources, or physical action.
\end{definition}

The pair $(I,O)$ completely characterizes the coupling between the system and
its environment.

Consequently, the environment never acts directly upon the internal state.
Every influence is mediated by an interface.

\begin{definition}[Admissible Interface]

An interface $(I,O)$ is called \emph{admissible} if every permitted exchange
preserves the invariant family

\[
\mathcal{J}
=
\{
J_1,
J_2,
\ldots,
J_n
\},
\]

where each

\[
J_i:\mathcal{S}\rightarrow\mathbb{R}
\]

is an invariant functional defined on the internal state space.

Equivalently,

\[
J_i(s)=J_i(I(e,s))
\]

for every admissible interaction.

\end{definition}

Thus admissibility is not defined by silence.

It is defined by invariant preservation.

A perfectly silent system may still destroy its own internal organization,
while a highly communicative system may remain completely admissible provided
its structural invariants remain unchanged.

\section{Coupling Operators}

Interactions are rarely binary.

Most systems permit certain forms of exchange while rejecting others.

This motivates the introduction of a coupling operator.

\begin{definition}[Coupling Operator]

A coupling operator is a bounded map

\[
\Gamma:
\mathcal{E}
\times
\mathcal{S}
\rightarrow
[0,1],
\]

assigning each possible interaction a coupling strength.

The extreme cases satisfy

\[
\Gamma=0
\]

for complete rejection and

\[
\Gamma=1
\]

for unrestricted acceptance.

Intermediate values describe partially admissible interactions.

\end{definition}

Notice that complete decoupling is neither necessary nor desirable.

If

\[
\Gamma\equiv0,
\]

then no external information may enter the system.

Learning becomes impossible.

Scientific collaboration disappears.

Repair cannot incorporate new evidence.

Likewise,

\[
\Gamma\equiv1
\]

is equally pathological.

Every perturbation propagates directly into internal dynamics.

Noise becomes indistinguishable from signal.

Autonomy collapses into environmental forcing.

The mathematically interesting regime therefore satisfies

\[
0 < \Gamma < 1,
\]

where selective interaction between the system and its environment becomes possible.

\begin{definition}[Selective Permeability]

An autonomous system is said to be \emph{selectively permeable} whenever its coupling operator is state dependent,

\[
\Gamma = \Gamma(e,s),
\]

so that admissibility depends jointly upon environmental structure and the system's internal configuration.

\end{definition}

Selective permeability is substantially more general than the architectural metaphor of a wall. Whereas a wall functions as a passive barrier that merely blocks or permits passage according to fixed physical constraints, an admissible interface acts as an active evaluative boundary. Rather than enforcing exclusion for its own sake, it computes whether proposed interactions satisfy the governing admissibility conditions and selectively permits only those transitions that preserve the system's long-term structural integrity. Its primary purpose is therefore not obstruction but the maintenance of coherence, stability, and persistence through principled evaluation.

\begin{proposition}

Every computationally active autonomous system possessing nontrivial learning dynamics necessarily admits a nonzero coupling operator.

\end{proposition}

\begin{proof}

Assume, for contradiction, that

\[
\Gamma \equiv 0.
\]

Then no information from the environment can enter the system. Consequently, the state evolution is completely determined by the internal dynamics,

\[
s_{t+1}=\Phi(s_t),
\]

independently of environmental observations. Since no external evidence can influence the system's evolution, Bayesian updating, statistical inference, parameter estimation, repair, and learning are all impossible. This contradicts the assumption that the system exhibits nontrivial adaptive dynamics. Therefore every adaptive computational system necessarily satisfies

\[
\Gamma \neq 0.
\]

\end{proof}

The proposition identifies nonvanishing coupling as a necessary condition for adaptive computation. A system with $\Gamma = 0$ is informationally closed and therefore incapable of incorporating evidence from its environment, whereas a system with $\Gamma > 0$ admits the possibility of observation, inference, learning, and repair. Selective permeability thus characterizes the regime in which an autonomous system can remain structurally coherent while continuing to adapt to changing external conditions.

This result establishes an important conceptual distinction between
adaptive computational systems and degenerate systems lacking a
distinguishable informational geometry. When \(\Gamma = 0\), all
configurations become informationally indistinguishable, preventing any
meaningful process of inference, learning, or repair. In contrast,
\(\Gamma \neq 0\) endows the system with a nontrivial geometry of
distinguishability, allowing evidence to discriminate among competing
hypotheses, parameters to be estimated, and adaptive corrections to
reduce reconstruction error. Thus, within the present framework,
nonvanishing distinguishability is not merely a technical assumption but
a necessary structural condition for admissible computation.

\chapter{Structural Entropy and Environmental Perturbation}

\section{Noise Is Not an Environmental Property}

Informal discussions frequently distinguish between environments that are
``peaceful'' and those that are ``noisy.'' Such language is useful
psychologically but lacks mathematical precision.

Within the present framework, noise is not an intrinsic property of the
environment.

Rather, it is a property of the interaction between the environment and the
state of the autonomous system.

Consequently identical external events may constitute severe perturbations
for one system while remaining dynamically negligible for another.

Noise is therefore relational rather than absolute.

\begin{definition}[Environmental Perturbation]

Let

\[
e\in\mathcal{E}
\]

be an environmental event.

The perturbation induced by $e$ upon state $s$ is

\[
P(e,s)
=
\Gamma(e,s)\,
\Delta(e,s),
\]

where

\[
\Gamma
\]

is the coupling operator introduced previously and

\[
\Delta
\]

measures the structural displacement produced by the interaction.

\end{definition}

Notice that an event may possess arbitrarily large magnitude while producing
arbitrarily small perturbation whenever

\[
\Gamma(e,s)\approx0.
\]

Likewise, even relatively minor events may produce significant structural
distortion whenever coupling approaches unity.

Thus disturbance cannot be identified solely with external intensity.

\section{Structural Entropy}

The cumulative effect of perturbations is quantified through a structural
entropy functional.

\begin{definition}[Structural Entropy]

Let

\[
\mathcal{H}:
\mathcal{S}
\rightarrow
\mathbb{R}_{\ge0}
\]

measure the degree of internal organizational uncertainty.

For every admissible evolution

\[
\Phi_t,
\]

the entropy increment is

\[
\Delta\mathcal{H}
=
\mathcal{H}(\Phi_t(s))
-
\mathcal{H}(s).
\]

\end{definition}

Unlike thermodynamic entropy, structural entropy need not measure disorder in
a physical sense.

Instead it quantifies degradation of the organizational relations required
for coherent computation.

Consequently highly dynamic systems may possess extremely low structural
entropy, whereas apparently motionless systems may accumulate structural
entropy rapidly.

The relevant quantity is organization rather than activity.

\section{Signal}

If perturbation is relational, then signal must also be defined
relationally.

\begin{definition}[Signal]

An interaction

\[
e\in\mathcal{E}
\]

is called a signal whenever

\[
\Delta\mathcal{H}\le0
\]

or

\[
\Delta\mathcal{J}\ge0,
\]

where

\[
\mathcal{J}
\]

denotes the invariant family introduced previously.

Signals therefore either reduce structural entropy or strengthen invariant
preservation.

\end{definition}

This definition immediately separates usefulness from novelty.

Novel information need not constitute signal.

Likewise, repeated observations may continue functioning as signal whenever
they reinforce organizational stability.

\section{Noise}

Noise is defined as the complementary class.

\begin{definition}[Noise]

An interaction is noise whenever

\[
\Delta\mathcal{H}>0
\]

while producing no compensating increase in invariant preservation,
computational capability, or repair potential.

\end{definition}

Noise therefore possesses no independent existence.

It is simply environmental interaction whose informational cost exceeds its
structural benefit.

\begin{remark}

This definition differs fundamentally from communication theory.

Classical information theory measures uncertainty in transmitted symbols.

The present framework measures degradation of autonomous organizational
structure.

The two notions coincide only under additional assumptions.

\end{remark}

\section{The Filtering Principle}

An admissible interface should therefore reject neither the largest nor the
most frequent perturbations.

Instead it minimizes the expected increase in structural entropy.

Formally,

\[
\Gamma^{*}
=
\operatorname*{arg\,min}_{\Gamma}
\mathbb{E}
\left[
\Delta\mathcal{H}
\right],
\]

subject to the constraints

\[
\mathcal{J}_i
=
\text{constant},
\]

for every preserved invariant.

This optimization criterion replaces numerous informal concepts including
``protecting one's peace,'' ``building walls,'' or ``remaining untouched.''

The objective is neither silence nor isolation.

The objective is entropy minimization under admissibility constraints.

\begin{theorem}[Entropy-Admissibility Principle]

Suppose an autonomous system remains computationally active while satisfying

\[
\Delta\mathcal{H}\le0
\]

for every admissible interaction.

Then every infinite admissible trajectory remains structurally stable.

\end{theorem}

\begin{proof}

Since structural entropy is bounded below by zero and is non-increasing along
every admissible trajectory,

\[
\{\mathcal{H}(s_n)\}
\]

forms a monotone bounded sequence.

Hence

\[
\lim_{n\rightarrow\infty}
\mathcal{H}(s_n)
\]

exists.

Because invariant preservation prevents catastrophic structural collapse,
the limiting trajectory remains admissible for all finite time.

Therefore long-term computational stability follows.

\end{proof}

The conclusion is notable.

A sanctuary is not characterized by the absence of disturbance.

Rather, it is characterized by the existence of interface operators that
prevent disturbance from accumulating into irreversible structural entropy.

\chapter{Repair Operators and Recursive Homeostasis}

\section{Home Is Not a Location}

The common identification of home with a geographical location, building, or
social group is mathematically inadequate.

Any location may lose its organizational coherence.

Conversely, a computational system may preserve internal identity while
undergoing continuous spatial relocation.

Consequently, home cannot be identified with position in Euclidean space.

Instead it must be regarded as an invariant subset of admissible
configurations.

\begin{definition}[Home]

Let

\[
\mathsf{Adm}\subseteq\mathcal{S}
\]

denote the admissible manifold.

The \emph{home} of an autonomous system is the maximal positively invariant
subset

\[
\mathcal{H}_{\mathrm{ome}}
\subseteq
\mathsf{Adm}
\]

such that

\[
\Phi_t(s)
\in
\mathcal{H}_{\mathrm{ome}}
\]

for every

\[
t\ge0
\]

whenever

\[
s\in
\mathcal{H}_{\mathrm{ome}}.
\]

\end{definition}

Thus home is not a coordinate.

Home is an invariant region of state space.

A room may contain home.

A laboratory may contain home.

A distributed computational network may contain home.

None of these physical realizations is itself home.

They merely instantiate admissible regions.

\section{Deviation from Admissibility}

No autonomous system remains permanently within the admissible manifold.

Unexpected observations, failures, interruptions, and incomplete information
continuously perturb the trajectory.

Consequently repair is unavoidable.

\begin{definition}[Defect]

The defect of a state

\[
s
\]

is measured by

\[
\delta(s)
=
d
\left(
s,
\mathsf{Adm}
\right),
\]

where

\[
d
\]

is any admissible metric or pseudometric compatible with the topology of
$\mathcal{S}$.

\end{definition}

States satisfying

\[
\delta(s)=0
\]

are admissible.

Positive values quantify the degree of structural displacement.

Importantly,

\[
\delta
\]

measures neither moral failure nor psychological distress.

It measures geometric distance from invariant organization.

\section{Repair Operators}

The central dynamical object of the theory is the repair operator.

\begin{definition}[Repair Operator]

A repair operator is a mapping

\[
R:
\mathcal{S}
\rightarrow
\mathcal{S}
\]

satisfying

\[
\delta(R(s))
\le
\delta(s)
\]

for every state

\[
s.
\]

If the inequality is strict whenever

\[
\delta(s)>0,
\]

the repair operator is called
\emph{strictly restorative}.

\end{definition}

Repair therefore need not return the system to its previous state.

Its only mathematical obligation is to decrease structural defect.

History is not reversed.

Admissibility is recovered.

\begin{remark}

Repair is therefore fundamentally different from restoration.

Restoration attempts to recreate the past.

Repair constructs an admissible future.

\end{remark}

\section{Recursive Repair}

Complex systems rarely admit complete repair in a single iteration.

Instead repair proceeds recursively.

Define

\[
R^{(0)}
=
\operatorname{id},
\]

and recursively

\[
R^{(n+1)}
=
R
\circ
R^{(n)}.
\]

The resulting sequence

\[
\{
R^{(n)}(s)
\}
\]

defines a repair trajectory.

\begin{definition}[Recursive Homeostasis]

An autonomous system exhibits recursive homeostasis whenever

\[
\lim_{n\rightarrow\infty}
\delta
\left(
R^{(n)}(s)
\right)
=
0
\]

for every admissible initial defect.

\end{definition}

Homeostasis therefore emerges as repeated approximation rather than static
equilibrium.

The system continually reconstructs itself.

\section{Home as an Attractor}

The previous definitions imply a stronger characterization.

\begin{theorem}[Home Attractor Theorem]

Suppose

\begin{enumerate}
\item the repair operator is continuous;
\item the admissible manifold is closed;
\item repair is strictly restorative outside
$\mathsf{Adm}$.
\end{enumerate}

Then

\[
\mathcal{H}_{\mathrm{ome}}
\]

is an attractor of the repair dynamics.

\end{theorem}

\begin{proof}

Since repair decreases defect monotonically,

\[
\delta
\left(
R^{(n+1)}(s)
\right)
<
\delta
\left(
R^{(n)}(s)
\right)
\]

whenever

\[
\delta>0.
\]

Thus

\[
\{
\delta_n
\}
\]

forms a bounded monotone sequence.

Therefore

\[
\delta_n
\rightarrow
L
\ge0.
\]

If

\[
L>0,
\]

strict restoration would imply another decrease, contradicting convergence.

Hence

\[
L=0.
\]

Since the admissible manifold is closed,

\[
R^{(n)}(s)
\rightarrow
\mathcal{H}_{\mathrm{ome}}.
\]

Therefore home is an attracting invariant set.

\end{proof}

The consequence is subtle but important.

Home is not where a system begins.

Home is where recursive repair converges.

The distinction removes any dependence upon geography, ancestry, architecture,
or narrative symbolism.

Home becomes a purely dynamical object arising from invariant-preserving
repair.

The remainder of this monograph develops the consequences of treating
autonomous cognition, scientific communities, biological organisms, and
distributed computational systems as recursive repair processes evolving
toward admissible attractors rather than static equilibrium states.

\chapter{The Geometry of Collaboration}

\section{Against the Myth of Isolation}

A persistent misconception in both popular psychology and technical culture is
that creativity is fundamentally an isolated phenomenon.

This misconception arises from observing the localization of computation while
neglecting the topology of information exchange.

Although the internal computation of an autonomous system occurs within its
own state space, the informational substrate upon which that computation
depends is inherently distributed.

No theorem proves itself.

No scientific discipline constructs its own language.

No programming language implements its own compiler without inherited
infrastructure.

Autonomy therefore cannot imply informational independence.

Instead, autonomy concerns the preservation of internal organizational
structure while participating in a distributed computational ecology.

\section{Knowledge Graphs}

Scientific activity is naturally represented by a graph.

\begin{definition}[Knowledge Graph]

A knowledge graph is a directed graph

\[
\mathcal{G}
=
(V,E),
\]

whose vertices represent autonomous cognitive systems and whose directed
edges represent admissible informational exchange.

An edge

\[
u\rightarrow v
\]

exists whenever information originating at

\[
u
\]

becomes structurally incorporated into the admissible state of

\[
v.
\]

\end{definition}

Edges need not correspond to direct communication.

Books, software, published proofs, observations, experimental apparatus,
historical traditions, and mathematical notation all induce admissible edges
within the graph.

Consequently scientific collaboration is substantially richer than
interpersonal communication.

\section{Indirect Coupling}

Many collaborative systems exhibit no direct interaction whatsoever.

Instead they communicate through persistent informational artifacts.

\begin{definition}[Indirect Coupling]

Two autonomous systems

\[
A
\]

and

\[
B
\]

are indirectly coupled whenever there exists an admissible sequence

\[
A
\rightarrow
v_1
\rightarrow
\cdots
\rightarrow
v_n
\rightarrow
B.
\]

The sequence need not preserve temporal simultaneity.

\end{definition}

Thus Euclid remains coupled to modern mathematics.

Newton remains coupled to orbital mechanics.

Shannon remains coupled to every communication protocol derived from
information theory.

The coupling persists despite the absence of temporal coexistence.

Knowledge propagates through graph topology rather than chronology.

\section{Structural Independence}

Informal discussions often confuse physical proximity with intellectual
dependence.

The two concepts are mathematically independent.

\begin{definition}[Structural Independence]

An autonomous system

\[
S
\]

is structurally independent whenever every admissible trajectory satisfies

\[
\frac{\partial\Phi}{\partial e}
\]

bounded independently of transient environmental perturbations.

\end{definition}

Structural independence therefore concerns bounded sensitivity.

It does not imply absence of communication.

Indeed, the most productive scientific systems frequently maintain large
numbers of admissible informational edges while remaining structurally
stable.

\section{Distributed Repair}

Repair itself is rarely local.

Suppose

\[
\mathcal{G}
=
(V,E)
\]

is a knowledge graph.

Associate to every vertex

\[
v
\]

a repair operator

\[
R_v.
\]

The global repair dynamics become

\[
R_{\mathcal{G}}
=
\prod_{v\in V}
R_v,
\]

where composition follows the partial order induced by graph topology.

Consequently local repair may propagate throughout the network without any
single vertex possessing complete knowledge.

Scientific communities therefore behave as distributed repair systems.

Each participant repairs only a fragment of the total structure.

The resulting global reconstruction exceeds the capability of every
individual contributor.

\section{Emergent Intelligence}

The previous observations motivate a more general definition.

\begin{definition}[Emergent Intelligence]

Let

\[
\mathcal{G}
=
(V,E)
\]

be a knowledge graph.

Emergent intelligence is the increase in repair capacity obtained through
graph composition:

\[
\mathcal{R}(\mathcal{G})
>
\sum_{v\in V}
\mathcal{R}(v),
\]

where

\[
\mathcal{R}
\]

denotes the repair functional.

\end{definition}

The inequality expresses genuine emergence.

The graph performs computational work that cannot be attributed to any
individual vertex.

This phenomenon appears throughout mathematics, scientific collaboration,
distributed software engineering, biological ecosystems, and neural systems.

\section{The Fallacy of the Solitary Genius}

The mythology of the solitary genius arises from observing only localized
computation.

It ignores the graph upon which that computation depends.

Every theorem presupposes notation.

Every notation presupposes language.

Every language presupposes a historical community.

Every computation presupposes inherited infrastructure.

Therefore no autonomous system is graph-independent.

The proper mathematical object is not the isolated individual but the
distributed repair graph within which individual cognition occupies a single
vertex.

Autonomy consists not in severing edges but in maintaining admissible edges
whose cumulative dynamics preserve structural invariants while enabling
recursive repair across the graph.

\chapter{Persistence, Continuation, and Identity}

\section{Identity Is Not Substance}

Classical metaphysical accounts frequently regard identity as an intrinsic
property of an object. Such approaches implicitly assume that persistence is
explained by the continued existence of an underlying substance.

The present framework adopts a fundamentally different viewpoint.

Identity is not primitive.

Identity is an emergent invariant generated by recursively admissible
continuation.

Consequently, persistence is not explained by substance.

Substance is explained by persistence.

\section{Continuation Operators}

Every autonomous system evolves through time.

Let

\[
\Phi_t:\mathcal{S}\rightarrow\mathcal{S}
\]

denote the temporal evolution operator introduced previously.

Not every temporal evolution preserves identity.

We therefore distinguish arbitrary evolution from admissible continuation.

\begin{definition}[Continuation Operator]

A continuation operator is a family

\[
C_t:\mathcal{S}\rightarrow\mathcal{S}
\]

satisfying

\[
C_0=\operatorname{id}
\]

and

\[
C_{t+s}
=
C_t
\circ
C_s
\]

while preserving the invariant family

\[
\mathcal{J}.
\]

\end{definition}

The semigroup property expresses temporal coherence.

Identity therefore depends not upon static resemblance but upon lawful
continuation.

\section{Persistence}

Persistence admits a quantitative characterization.

\begin{definition}[Persistence Functional]

Let

\[
P:\mathcal{S}\rightarrow\mathbb{R}_{\ge0}
\]

be defined by

\[
P(s)
=
\int_0^\infty
w(t)
\,\chi_{\mathsf{Adm}}
(C_t(s))
\,dt,
\]

where

\[
w(t)
\]

is a nonnegative weighting function and

\[
\chi_{\mathsf{Adm}}
\]

is the characteristic function of the admissible manifold.

\end{definition}

Persistence therefore measures the cumulative duration over which admissible
continuation remains possible.

It is neither age nor longevity.

Rather, it measures structural survivability.

\section{Identity Through Repair}

Suppose perturbation produces

\[
s'
\notin
\mathsf{Adm}.
\]

Repair generates

\[
R(s').
\]

The repaired state generally differs from the original.

Nevertheless identity may remain preserved.

\begin{definition}[Identity Equivalence]

Two states

\[
s_1,s_2
\]

are identity equivalent whenever there exists a finite sequence

\[
s_1
\rightarrow
R_1
\rightarrow
R_2
\rightarrow
\cdots
\rightarrow
R_n
\rightarrow
s_2
\]

such that every intermediate state belongs to the admissible manifold.

We write

\[
s_1
\sim
s_2.
\]

\end{definition}

Identity therefore becomes an equivalence relation generated by admissible
repair rather than perfect state equality.

This distinction is essential.

Exact equality almost never survives long temporal evolution.

Admissible continuation frequently does.

\section{The Failure of Static Ontologies}

Suppose identity required

\[
s_t=s_0
\]

for every

\[
t.
\]

Then every nontrivial evolution would immediately destroy identity.

Biological development would constitute replacement.

Learning would constitute replacement.

Scientific progress would constitute replacement.

Computation itself would constitute replacement.

This conclusion is absurd.

Static identity therefore cannot provide an adequate ontology for evolving
systems.

\section{The Continuation Principle}

The previous constructions motivate the central principle.

\begin{theorem}[Continuation Principle]

Let

\[
\{C_t\}_{t\ge0}
\]

be a family of admissible continuation operators.

Suppose

\[
R
\]

restores every finite defect while preserving invariant equivalence.

Then identity is preserved along every recursively repaired continuation
trajectory.

\end{theorem}

\begin{proof}

Every admissible continuation preserves the invariant family

\[
\mathcal{J}.
\]

Whenever perturbation produces temporary inadmissibility,

\[
R
\]

returns the trajectory to the admissible manifold without destroying
equivalence.

Since composition of admissible continuation with admissible repair remains
admissible,

\[
C_t
\circ
R
\]

preserves the equivalence class generated by

\[
\sim.
\]

Repeated application therefore constructs an infinite admissible trajectory
contained entirely within a single identity class.

\end{proof}

\section{Identity as a Fixed Structure}

The previous theorem reverses a common intuition.

Identity is not the object that persists through change.

Identity is the invariant generated by lawful change.

Objects are transient realizations.

Identity is the stable geometry of their admissible continuation.

Accordingly, questions concerning the persistence of persons, institutions,
languages, scientific theories, computational systems, and biological
organisms reduce to a single mathematical problem:

\begin{quote}
Does there exist a recursively admissible continuation operator preserving
the relevant invariant family?
\end{quote}

The remainder of this work demonstrates that this question admits a common
answer across seemingly unrelated domains, thereby replacing diverse
metaphysical accounts of persistence with a unified theory of recursive
continuation.

\chapter{Observers, Interfaces, and Representation}

\section{The Impossibility of Direct Observation}

Every observational theory begins with an implicit assumption regarding the
relationship between an observer and the external world.

The simplest assumption is direct realism:

\[
\mathcal{O}
=
\mathcal{W},
\]

where observation is identified with reality itself.

Such an assumption is untenable.

No observer possesses unrestricted access to the complete external state of
its environment.

Every observation is mediated by finite sensors, finite computation,
measurement uncertainty, temporal delay, and representational constraints.

Consequently all observation is indirect.

The observer never receives reality.

The observer receives an interface.

\section{Interface States}

Let

\[
\mathcal{W}
\]

denote the external world.

Let

\[
\mathcal{I}
\]

denote the interface accessible to an observer.

The observer itself possesses internal state space

\[
\mathcal{S}.
\]

Observation therefore decomposes into two mappings

\[
\Pi:
\mathcal{W}
\rightarrow
\mathcal{I},
\]

and

\[
\mathcal{O}:
\mathcal{I}
\rightarrow
\mathcal{S}.
\]

The first mapping represents physical measurement.

The second represents interpretation.

Reality therefore satisfies the composition

\[
\mathcal{W}
\longrightarrow
\mathcal{I}
\longrightarrow
\mathcal{S},
\]

rather than

\[
\mathcal{W}
\longrightarrow
\mathcal{S}.
\]

The distinction is fundamental.

Errors introduced by the interface cannot be distinguished from properties of
the external world unless additional structural assumptions are available.

\section{Representation}

The observer does not manipulate reality directly.

Instead it manipulates internal representations.

\begin{definition}[Representation]

A representation is a mapping

\[
\rho:
\mathcal{I}
\rightarrow
\mathcal{R},
\]

where

\[
\mathcal{R}
\]

is an internal representational manifold.

Reasoning, planning, prediction, memory, and decision-making occur within

\[
\mathcal{R},
\]

not within

\[
\mathcal{W}.
\]

\end{definition}

Representations therefore possess operational reality even when they are
imperfect descriptions of the external world.

Scientific models, programming languages, maps, mathematical notation,
neural activity, and symbolic logic all constitute representational systems.

\section{Information Loss}

Every interface necessarily discards information.

\begin{definition}[Projection]

An observational interface is a projection

\[
\Pi:
\mathcal{W}
\rightarrow
\mathcal{I}
\]

whenever

\[
\dim(\mathcal{I})
<
\dim(\mathcal{W}),
\]

or, more generally, whenever distinct external configurations possess
identical interface states.

\end{definition}

Consequently

\[
\Pi
\]

cannot generally admit an inverse.

Different external realities may generate identical observations.

Observation therefore determines an equivalence class rather than a unique
world state.

\section{Observational Equivalence}

\begin{definition}[Observational Equivalence]

Two external configurations

\[
w_1,w_2
\in
\mathcal{W}
\]

are observationally equivalent whenever

\[
\Pi(w_1)
=
\Pi(w_2).
\]

We write

\[
w_1
\equiv_{\Pi}
w_2.
\]

\end{definition}

The observer cannot distinguish members of the same equivalence class using
the given interface alone.

Additional measurements, alternative interfaces, or prior structural
knowledge become necessary.

Thus ignorance is not always a failure of reasoning.

It may simply reflect insufficient observational resolution.

\section{Reconstruction}

Scientific inference attempts to reconstruct the external world from
interface observations.

This inversion is generally ill-posed.

\begin{definition}[Reconstruction Operator]

A reconstruction operator is a mapping

\[
\mathcal{R}:
\mathcal{I}
\rightarrow
2^{\mathcal{W}},
\]

assigning every interface state the admissible family of compatible external
configurations.

\end{definition}

Notice that

\[
\mathcal{R}(i)
\]

is typically not a singleton.

Observation constrains reality without uniquely determining it.

Inference therefore consists of progressively shrinking admissible
equivalence classes.

\section{Representation Defect}

The quality of a representation depends upon its ability to preserve relevant
structure.

\begin{definition}[Representation Defect]

Let

\[
\rho
\]

be an internal representation.

The representation defect is the functional

\[
\Delta_{\rho}
=
d
\!\left(
\Pi(\mathcal{W}),
\rho(\mathcal{I})
\right),
\]

where

\[
d
\]

is any admissible structural metric.

\end{definition}

Perfect representation satisfies

\[
\Delta_{\rho}=0.
\]

In practice,

\[
\Delta_{\rho}>0
\]

for every finite observer.

Scientific progress therefore consists not in eliminating representation
defect but in recursively reducing it.

\section{The Representation Principle}

The preceding constructions yield a general theorem.

\begin{theorem}[Representation Principle]

No finite observer possesses direct access to external reality.

Every autonomous cognitive system operates exclusively upon internal
representations generated through admissible observational interfaces.

Consequently every act of reasoning is necessarily reconstructive rather than
directly perceptual.

\end{theorem}

\begin{proof}

Observation is mediated by

\[
\Pi:
\mathcal{W}
\rightarrow
\mathcal{I}.
\]

Reasoning operates only upon

\[
\rho(\mathcal{I}).
\]

Since

\[
\Pi
\]

generally lacks an inverse, direct recovery of

\[
\mathcal{W}
\]

is impossible.

Therefore cognition necessarily proceeds by constructing admissible
reconstructions whose representation defect decreases through recursive
repair.

\end{proof}

The consequences extend well beyond epistemology.

Scientific theories, software systems, biological perception, artificial
intelligence, language, memory, and mathematical abstraction all become
instances of a common representational architecture.

Reality itself is not manipulated directly.

Only admissible representations evolve under recursive continuation while
remaining constrained by observational interfaces.

\chapter{Distinction Geometry and the Topology of Information}

\section{Information Is Not Primitive}

Traditional information theory begins by assuming the existence of symbols,
messages, or random variables.

The present framework begins earlier.

Prior to every message, prior to every probability distribution, and prior to
every symbolic language lies a more fundamental object:

the distinction.

Without distinguishable configurations, no observation, measurement,
communication, or computation can exist.

Consequently information is not primitive.

Information is induced by distinguishability.

\section{The Distinction Space}

Let

\[
\mathcal{X}
\]

be a configuration space.

We define an equivalence relation

\[
\sim
\]

representing observational indistinguishability.

The corresponding quotient space

\[
\mathcal{D}
=
\mathcal{X}/\sim
\]

is called the distinction space.

Its elements are equivalence classes

\[
[x]
=
\{
y\in\mathcal{X}
:
y\sim x
\}.
\]

Every observation identifies an element of

\[
\mathcal{D},
\]

not necessarily an element of

\[
\mathcal{X}.
\]

Thus cognition operates upon distinction classes rather than absolute states.

\section{Distinction Fields}

Distinctions are rarely isolated.

Instead they organize into continuous structures.

\begin{definition}[Distinction Field]

A distinction field is a mapping

\[
\Delta:
\mathcal{X}
\rightarrow
\mathbb{R}_{\ge0}
\]

whose value measures the local distinguishability surrounding a point.

Large values correspond to regions of high observational separation.

Small values indicate observational collapse.

\end{definition}

Unlike probability densities, distinction fields describe geometric
resolution rather than uncertainty.

They measure how easily neighboring configurations may be separated by an
admissible observer.

\section{Observational Collapse}

Finite interfaces cannot preserve arbitrary distinctions.

Projection necessarily collapses certain regions.

\begin{definition}[Collapse Region]

A subset

\[
U
\subseteq
\mathcal{X}
\]

is called a collapse region whenever

\[
\Pi(U)
=
\{i\}
\]

for some interface state

\[
i
\in
\mathcal{I}.
\]

\end{definition}

Within a collapse region every configuration becomes observationally
equivalent.

No internal computation can distinguish them using the given interface.

Collapse therefore reflects finite observational resolution rather than loss
of objective structure.

\section{Distinction Entropy}

The richness of an observational system depends upon the diversity of
distinctions it preserves.

\begin{definition}[Distinction Entropy]

Let

\[
\mu
\]

be a measure on

\[
\mathcal{D}.
\]

The distinction entropy is the functional

\[
\mathfrak{D}
=
\int_{\mathcal{D}}
\Delta(x)
\,d\mu(x).
\]

\end{definition}

This quantity measures the total distinguishable structure available to an
observer.

Unlike Shannon entropy, distinction entropy does not quantify uncertainty.

Instead it quantifies representational resolution.

Two observers possessing identical Shannon information may possess different
distinction entropies if one preserves finer geometric structure than the
other.

\section{Repair of Distinctions}

Repair does not merely restore admissibility.

It restores distinguishability.

Suppose perturbation merges two previously distinguishable regions.

The repair operator acts by increasing local distinction density.

Formally,

\[
R:
\Delta
\longrightarrow
\Delta',
\]

where

\[
\Delta'
\ge
\Delta
\]

whenever structural distinctions have been degraded.

Thus repair reconstructs the geometry upon which subsequent computation
depends.

\section{Distinction Curvature}

Local distinguishability may vary continuously.

This motivates a differential description.

\begin{definition}[Distinction Curvature]

The distinction curvature is the tensor

\[
\mathcal{K}
=
\nabla^2
\Delta,
\]

whenever the distinction field is twice differentiable.

\end{definition}

Regions of positive curvature correspond to rapidly increasing observational
resolution.

Regions of negative curvature correspond to distinction collapse.

Flat regions represent observational stability.

Thus the geometry of cognition may be analyzed using the same differential
tools employed throughout modern geometry.

\section{The Distinction Principle}

The preceding constructions yield a general characterization of information.

\begin{theorem}[Distinction Principle]

Every informational process factors through distinction geometry.

Specifically,

\[
\text{Information}
=
\text{Observation}
\circ
\text{Distinction}.
\]

Consequently, whenever distinguishability vanishes, information necessarily
vanishes regardless of computational complexity.

Conversely, increasing computational power cannot recover distinctions that
have already been destroyed by projection.

\end{theorem}

\begin{proof}

Observation first partitions the external configuration space into
equivalence classes.

Internal computation operates only upon these classes.

Therefore every subsequent informational process is defined upon the
distinction space

\[
\mathcal{D},
\]

rather than upon the original configuration space.

Hence distinguishability is logically prior to information.

\end{proof}

The theory developed in this chapter reverses the conventional hierarchy.

Classical information theory regards distinctions as consequences of encoded
messages.

The present framework regards encoded messages as consequences of
pre-existing distinguishable structure.

Information therefore emerges from geometry rather than geometry emerging from
information.

This inversion provides a common mathematical foundation for perception,
scientific inference, communication theory, artificial intelligence, and
recursive repair.

\chapter{Constraint Geometry and Reachability}

\section{Constraint Precedes Motion}

Classical mechanics begins by specifying trajectories.

Optimization begins by specifying objective functions.

Control theory begins by specifying admissible inputs.

The present framework begins earlier.

Before any trajectory exists, before any optimization may be attempted, and
before any computation may occur, there exists a geometry of possibility.

That geometry is determined entirely by constraints.

Motion therefore does not define constraints.

Constraints define motion.

\section{Constraint Manifolds}

Let

\[
\mathcal{S}
\]

denote the complete configuration space.

Only a subset corresponds to physically, computationally, or logically
realizable configurations.

\begin{definition}[Constraint Manifold]

A constraint manifold is a subset

\[
\mathcal{C}
\subseteq
\mathcal{S}
\]

defined by a family of admissibility relations

\[
F_i(s)=0,
\]

together with inequality constraints

\[
G_j(s)\ge0.
\]

The manifold consists precisely of those configurations satisfying every
constraint simultaneously.

\end{definition}

Every autonomous system evolves entirely within

\[
\mathcal{C},
\]

except during transient perturbations subsequently corrected by repair.

Thus constraints do not oppose dynamics.

They define the domain upon which dynamics exist.

\section{Reachability}

Not every admissible configuration can necessarily evolve into every other.

This observation motivates the central geometric object.

\begin{definition}[Reachability]

Given two admissible states

\[
s_0,s_1
\in
\mathcal{C},
\]

we say that

\[
s_1
\]

is reachable from

\[
s_0
\]

whenever there exists a continuous admissible trajectory

\[
\gamma:[0,1]\rightarrow\mathcal{C}
\]

such that

\[
\gamma(0)=s_0,
\qquad
\gamma(1)=s_1.
\]

\end{definition}

Reachability is therefore a topological property rather than merely a
kinematic one.

The existence of a destination does not imply the existence of an admissible
path.

\section{Reachability Components}

Reachability induces an equivalence relation upon admissible states.

\begin{definition}[Reachability Component]

The reachability component of

\[
s
\]

is

\[
\mathcal{R}(s)
=
\{
x\in\mathcal{C}
:
x
\leftrightsquigarrow
s
\},
\]

where

\[
\leftrightsquigarrow
\]

denotes mutual reachability.

\end{definition}

Thus

\[
\mathcal{C}
=
\bigsqcup_i
\mathcal{R}_i,
\]

decomposes into disjoint connected components.

Each component represents an autonomous region of possible evolution.

Repair may move trajectories within a component.

Structural innovation may enlarge a component.

Constraint modification may merge components.

\section{Boundary Geometry}

The boundary of a reachability component possesses particular significance.

\begin{definition}[Reachability Boundary]

The reachability boundary is

\[
\partial\mathcal{R}
=
\overline{\mathcal{R}}
\setminus
\operatorname{int}(\mathcal{R}).
\]

\end{definition}

Near the boundary, arbitrarily small perturbations may destroy admissibility.

Consequently boundary states exhibit maximal structural sensitivity.

Conversely, interior states possess robustness under perturbation.

This distinction suggests that stable autonomous systems naturally evolve
toward interior regions rather than remaining arbitrarily close to constraint
boundaries.

\section{Constraint Curvature}

Constraints possess local geometry.

Suppose

\[
\mathcal{C}
\]

is a differentiable manifold.

The second fundamental form determines local bending relative to the ambient
configuration space.

Regions of large curvature correspond to rapidly changing admissibility.

Regions of small curvature admit comparatively stable continuation.

Thus computational difficulty frequently arises not from dimensionality but
from curvature.

Highly curved constraint surfaces require continual corrective action to
remain admissible.

\section{Constraint Energy}

Deviation from admissibility admits a scalar potential.

\begin{definition}[Constraint Energy]

The constraint energy is a functional

\[
E_C:
\mathcal{S}
\rightarrow
\mathbb{R}_{\ge0}
\]

satisfying

\[
E_C(s)=0
\]

if and only if

\[
s\in\mathcal{C},
\]

and

\[
E_C(s)>0
\]

otherwise.

\end{definition}

Repair operators may therefore be interpreted as gradient-like flows

\[
\frac{ds}{dt}
=
-
\nabla
E_C(s),
\]

whenever differentiability assumptions hold.

Repair becomes geometric descent toward admissibility.

\section{The Reachability Principle}

The previous constructions yield the central theorem of this chapter.

\begin{theorem}[Reachability Principle]

The behavior of an autonomous system is determined less by the collection of
states it may occupy than by the topology of admissible paths connecting
those states.

\end{theorem}

\begin{proof}

State descriptions alone specify isolated configurations.

Autonomous evolution requires continuous admissible trajectories.

Consequently every achievable computation, biological process, scientific
discovery, or cognitive transition depends upon the existence of paths within
the constraint manifold.

Removing paths while preserving states destroys behavior without altering the
underlying state space.

Hence reachability constitutes the primary dynamical invariant.

\end{proof}

This theorem shifts emphasis from configuration to navigation.

The relevant scientific question is no longer

\begin{quote}
``What states exist?''
\end{quote}

but instead

\begin{quote}
``Which admissible states can be reached, repaired, continued, and preserved
without violating structural constraints?''
\end{quote}

Accordingly, dynamics becomes the geometry of admissible navigation.

Constraints determine topology.

Topology determines reachability.

Reachability determines computation.

Computation determines the observable evolution of autonomous systems.

The remainder of this work develops this hierarchy into a unified theory of
physical, biological, cognitive, and computational organization.

\chapter{The Geometry of Constraint Propagation}

\section{Local Constraints and Global Structure}

A fundamental feature of complex systems is that global organization need not
be imposed globally.

Instead, large-scale coherence frequently emerges from purely local
constraints.

This phenomenon appears throughout mathematics and science.

Local differential equations generate global trajectories.

Local compatibility conditions generate smooth manifolds.

Local rewriting rules generate formal languages.

Local physical interactions generate biological organisms.

The common mathematical principle is therefore not global construction but
constraint propagation.

\section{Constraint Networks}

Individual constraints rarely exist in isolation.

Instead they form interconnected systems.

\begin{definition}[Constraint Network]

Let

\[
V=\{C_1,C_2,\ldots,C_n\}
\]

be a finite or countable family of constraints.

A constraint network is a graph

\[
\mathcal{N}=(V,E),
\]

where

\[
(C_i,C_j)\in E
\]

whenever satisfaction of one constraint influences the admissibility of the
other.

\end{definition}

The graph need not be complete.

Indeed, sparse constraint graphs frequently exhibit greater computational
efficiency while preserving equivalent admissible solution sets.

\section{Constraint Compatibility}

Not every family of constraints admits simultaneous realization.

\begin{definition}[Compatibility]

A collection of constraints

\[
\{C_i\}
\]

is compatible whenever

\[
\bigcap_i C_i
\neq
\varnothing.
\]

Otherwise the collection is incompatible.

\end{definition}

Compatibility therefore represents the existence of at least one admissible
configuration satisfying every constraint simultaneously.

Incompatibility implies structural inconsistency rather than computational
difficulty.

No algorithm may realize an impossible configuration.

\section{Propagation Operators}

Information acquired locally should modify neighboring constraints.

\begin{definition}[Propagation Operator]

A propagation operator

\[
P:
\mathcal{N}
\rightarrow
\mathcal{N}
\]

updates the admissible domains associated with every constraint by
communicating local information throughout the constraint graph.

\end{definition}

Propagation differs fundamentally from search.

Propagation removes impossible configurations.

Search chooses among remaining possibilities.

Thus propagation reduces complexity before decision.

\section{Monotonicity}

Constraint propagation should never enlarge the admissible search space.

\begin{proposition}

Suppose

\[
D_i^{(k)}
\]

denotes the admissible domain associated with constraint

\[
C_i
\]

after

\[
k
\]

propagation steps.

Then

\[
D_i^{(k+1)}
\subseteq
D_i^{(k)}
\]

for every

\[
i.
\]

\end{proposition}

\begin{proof}

Each propagation step removes assignments inconsistent with neighboring
constraints.

No previously eliminated configuration can become admissible without
modifying the constraint network itself.

Hence admissible domains decrease monotonically.

\end{proof}

Constraint propagation therefore possesses an intrinsic directionality.

Knowledge accumulates by eliminating impossibilities.

\section{Fixed Points}

Propagation eventually stabilizes.

\begin{definition}[Propagation Fixed Point]

A constraint network reaches a propagation fixed point whenever

\[
P(\mathcal{N})
=
\mathcal{N}.
\]

Equivalently,

\[
D_i^{(k+1)}
=
D_i^{(k)}
\]

for every constraint.

\end{definition}

At a fixed point, every remaining admissible configuration is locally
consistent with every propagated constraint.

Further propagation alone produces no additional information.

Subsequent progress requires either new observations or explicit search.

\section{Constraint Repair}

Perturbations may invalidate previously established consistency.

Repair therefore acts upon the network itself.

\begin{definition}[Constraint Repair]

A repair operator

\[
R_{\mathcal{N}}
:
\mathcal{N}
\rightarrow
\mathcal{N}
\]

is admissible whenever it restores compatibility while minimizing structural
modification of the original network.

\end{definition}

Repair therefore preserves as much existing organization as possible.

It does not reconstruct the network from first principles.

Instead it performs minimal admissible correction.

\section{Distributed Consistency}

The interaction between propagation and repair produces an important
principle.

\begin{theorem}[Distributed Consistency Principle]

Suppose every local constraint satisfies compatibility with each of its
neighbors after propagation.

Further suppose repair restores every finite incompatibility introduced by
external perturbation.

Then global consistency emerges as the least fixed point of alternating
propagation and repair.

\end{theorem}

\begin{proof}

Propagation monotonically reduces admissible domains while preserving every
globally realizable configuration.

Repair restores compatibility whenever local inconsistencies arise.

Since admissible domains form descending chains bounded below by the empty
set, propagation converges.

Repair prevents convergence to incompatible states.

Consequently the alternating composition

\[
P
\circ
R_{\mathcal{N}}
\]

approaches the least compatible fixed point containing all admissible
solutions.

\end{proof}

This theorem illustrates a recurring theme throughout the present work.

Global order need not be explicitly imposed.

Instead, sufficiently rich systems construct coherent large-scale structure
through recursive interaction between local admissibility, propagation,
repair, and continuation.

Accordingly, complexity is not generated by centralized control.

It emerges from recursively maintained consistency across distributed
constraint networks.

\chapter{Sheaf-Theoretic Reconstruction and Global Consistency}

\section{The Local Nature of Knowledge}

Every sufficiently complex observational process is fundamentally local.

No observer measures an entire physical system simultaneously.

No scientific discipline possesses complete knowledge.

No computational process evaluates every possible state at once.

Instead, knowledge is acquired through overlapping local observations whose
mutual compatibility determines the existence of a coherent global model.

Consequently, the mathematical problem is not merely one of observation.

It is one of reconstruction.

\section{Observation Covers}

Let

\[
\mathcal{X}
\]

denote the external configuration manifold.

Suppose

\[
\mathcal{U}
=
\{
U_i
\}_{i\in I}
\]

is an open cover satisfying

\[
\mathcal{X}
=
\bigcup_{i\in I}
U_i.
\]

Each

\[
U_i
\]

represents a region accessible to an individual observer, experiment,
measurement device, computational subsystem, or scientific discipline.

No single element of the cover need contain complete information.

Completeness arises only through compatibility across intersections.

\section{Local Sections}

Observations assign internal descriptions to local regions.

\begin{definition}[Local Section]

A local section over

\[
U_i
\]

is a mapping

\[
s_i
:
U_i
\rightarrow
\mathcal{R},
\]

where

\[
\mathcal{R}
\]

denotes the representational space introduced previously.

\end{definition}

The collection

\[
\{
s_i
\}
\]

therefore represents distributed partial knowledge rather than a complete
description of the external manifold.

\section{Compatibility}

Independent observations need not contradict one another.

Their compatibility is determined on overlaps.

\begin{definition}[Compatibility Condition]

Two local sections

\[
s_i,
s_j
\]

are compatible whenever

\[
s_i|_{U_i\cap U_j}
=
s_j|_{U_i\cap U_j}.
\]

\end{definition}

Thus agreement is required only where observations overlap.

Outside the intersection, the sections remain independent.

Compatibility therefore constitutes a local condition.

\section{The Gluing Principle}

The central mathematical question becomes the following.

Given compatible local observations,

does there exist a unique global reconstruction?

\begin{definition}[Global Section]

A global section is a mapping

\[
S
:
\mathcal{X}
\rightarrow
\mathcal{R}
\]

whose restriction satisfies

\[
S|_{U_i}
=
s_i
\]

for every member of the cover.

\end{definition}

Whenever such an

\[
S
\]

exists, the local observations admit coherent global interpretation.

\section{Reconstruction Defect}

Failure of reconstruction admits quantitative measurement.

\begin{definition}[Reconstruction Defect]

The reconstruction defect associated with an observation cover is

\[
\Delta_{\mathrm{rec}}
=
\sum_{i,j}
d
\!\left(
s_i|_{U_i\cap U_j},
s_j|_{U_i\cap U_j}
\right),
\]

where

\[
d
\]

is an admissible metric on representational space.

\end{definition}

Perfect reconstruction satisfies

\[
\Delta_{\mathrm{rec}}
=
0.
\]

Positive values indicate disagreement between neighboring observations.

Thus reconstruction becomes an optimization problem rather than a binary
decision.

\section{Repair by Gluing}

Repair now acquires a geometric interpretation.

Rather than modifying isolated observations, repair minimizes overlap
inconsistency.

Formally,

\[
R
:
\{
s_i
\}
\longrightarrow
\{
s_i'
\},
\]

where

\[
\Delta_{\mathrm{rec}}
(s_i')
<
\Delta_{\mathrm{rec}}
(s_i).
\]

Repair therefore increases global coherence by improving local compatibility.

Notice that no observer requires complete knowledge.

Each observer need only improve agreement with neighboring sections.

\section{Distributed Truth}

The preceding construction motivates an alternative conception of scientific
truth.

Truth is not the possession of an individual observer.

Nor is it the consequence of a privileged viewpoint.

Instead, truth emerges whenever distributed local observations admit a stable
global reconstruction.

Scientific collaboration therefore becomes an iterative gluing process.

Each experiment extends one section.

Each proof enlarges another.

Each computation refines a third.

Global understanding appears only through recursive compatibility.

\section{The Local-to-Global Principle}

The previous constructions culminate in the principal theorem.

\begin{theorem}[Local-to-Global Reconstruction]

Suppose

\[
\mathcal{U}
=
\{
U_i
\}
\]

is an observation cover.

Assume

\begin{enumerate}
\item every local section is admissible,
\item neighboring sections satisfy compatibility,
\item recursive repair decreases reconstruction defect.
\end{enumerate}

Then repeated repair converges toward a global section whenever one exists.

\end{theorem}

\begin{proof}

Compatibility defines local consistency on every overlap.

Repair monotonically decreases

\[
\Delta_{\mathrm{rec}}.
\]

Since

\[
\Delta_{\mathrm{rec}}
\ge
0,
\]

the reconstruction sequence possesses a lower bound.

Under continuity of repair and compactness of the admissible representational
space, convergence follows.

The limiting family of sections therefore defines a global reconstruction
compatible with every local observation.

\end{proof}

This theorem establishes a profound inversion of classical epistemology.

Knowledge is not assembled by beginning with a complete global model and
analyzing its parts.

Rather, coherent global structure emerges from recursively repairing
relationships among overlapping local observations.

Accordingly, mathematics, scientific inquiry, distributed computation,
biological cognition, and collaborative intelligence become instances of a
single geometric process:

\[
\boxed{
\text{Local Observation}
\rightarrow
\text{Compatibility}
\rightarrow
\text{Repair}
\rightarrow
\text{Global Reconstruction}.
}
\]

\chapter{Functorial Observation and the Category of Autonomous Systems}

\section{Beyond State Spaces}

Previous chapters developed autonomous systems as state spaces equipped with
admissibility, repair, continuation, and interface operators.

Although this formulation is sufficient for many applications, it conceals an
important structural fact.

The essential object is not the state itself.

The essential object is the collection of transformations that preserve
admissibility.

Consequently, the natural mathematical language is category theory.

\section{The Category of Autonomous Systems}

\begin{definition}[Autonomous Category]

Define the category

\[
\mathbf{Aut}
\]

as follows.

Objects are autonomous systems

\[
(\mathcal{S},\Phi,\mathcal{J}),
\]

where

\begin{enumerate}
\item $\mathcal{S}$ is a configuration space,
\item $\Phi$ denotes admissible evolution,
\item $\mathcal{J}$ is the invariant family.
\end{enumerate}

Morphisms

\[
f:A\rightarrow B
\]

are admissibility-preserving transformations satisfying

\[
f(\mathsf{Adm}_A)
\subseteq
\mathsf{Adm}_B.
\]

\end{definition}

Composition is ordinary function composition,

\[
g\circ f,
\]

while identity morphisms are

\[
\operatorname{id}_{\mathcal S}.
\]

Thus autonomous systems naturally form a category.

\section{Observation as a Functor}

Observation does not preserve every property of reality.

Instead it preserves only admissible structure.

This observation motivates a functorial description.

\begin{definition}[Observation Functor]

The observation process is a functor

\[
\mathcal{O}
:
\mathbf{Aut}
\longrightarrow
\mathbf{Rep},
\]

where

\[
\mathbf{Rep}
\]

denotes the category of internal representations.

\end{definition}

The functor satisfies

\[
\mathcal{O}
(\operatorname{id})
=
\operatorname{id},
\]

and

\[
\mathcal{O}(g\circ f)
=
\mathcal{O}(g)
\circ
\mathcal{O}(f).
\]

Observation therefore preserves compositional structure while generally
discarding geometric detail.

\section{Forgetful Processes}

Every observation loses information.

Categorically, this is expressed by forgetful functors.

\begin{definition}[Forgetful Functor]

A forgetful functor

\[
U
:
\mathbf{Aut}
\rightarrow
\mathbf{Set}
\]

maps every autonomous system to its underlying collection of states while
discarding repair operators, admissibility relations, continuation dynamics,
and invariant structure.

\end{definition}

Many traditional scientific descriptions are implicitly forgetful.

They retain observable variables while discarding the morphisms responsible
for long-term organization.

The present framework instead treats those morphisms as primary.

\section{Natural Transformations}

Suppose two observers construct different internal representations.

Each defines a functor

\[
\mathcal{O}_1,
\qquad
\mathcal{O}_2.
\]

Agreement between observers should not require identical internal
representations.

Instead it requires structural compatibility.

\begin{definition}[Natural Interpretation]

A natural transformation

\[
\eta:
\mathcal{O}_1
\Rightarrow
\mathcal{O}_2
\]

is called an interpretation whenever every admissible morphism satisfies the
commutative diagram

\[
\eta_B
\circ
\mathcal{O}_1(f)
=
\mathcal{O}_2(f)
\circ
\eta_A.
\]

\end{definition}

Interpretation therefore preserves relationships rather than individual
descriptions.

Translation becomes structure-preserving rather than symbol-preserving.

\section{Repair as an Endofunctor}

Repair operates internally upon autonomous systems.

Consequently it defines an endofunctor.

\begin{definition}[Repair Endofunctor]

The repair operator determines an endofunctor

\[
\mathcal{R}
:
\mathbf{Aut}
\rightarrow
\mathbf{Aut},
\]

acting upon objects by decreasing structural defect while preserving
admissibility.

Morphisms are mapped to repaired morphisms compatible with repaired state
spaces.

\end{definition}

Repeated repair therefore corresponds to repeated application of

\[
\mathcal{R}.
\]

The recursive sequence

\[
\mathcal{R},
\quad
\mathcal{R}^2,
\quad
\mathcal{R}^3,
\ldots
\]

defines the categorical analogue of recursive continuation.

\section{Equivalence of Autonomous Systems}

Two systems need not possess identical internal realizations to exhibit
identical structural organization.

\begin{definition}[Structural Equivalence]

Autonomous systems

\[
A
\]

and

\[
B
\]

are structurally equivalent whenever there exists an equivalence of
categories preserving

\begin{enumerate}
\item admissibility,
\item repair,
\item continuation,
\item invariant preservation.
\end{enumerate}

\end{definition}

Thus identical implementation is unnecessary.

Only preservation of organizational structure is required.

This observation explains why identical computational principles frequently
appear in biological evolution, distributed software systems, mathematical
reasoning, and scientific collaboration despite radically different physical
substrates.

\section{The Functorial Principle}

The previous constructions culminate in the following theorem.

\begin{theorem}[Functorial Principle]

Every admissible observational process factors through a functor between
categories of autonomous systems.

Consequently, scientific explanation consists not in cataloguing objects but
in identifying functorially preserved structure.

\end{theorem}

\begin{proof}

Observation preserves identity morphisms and composition.

Therefore it defines a functor.

Repair defines an endofunctor preserving admissibility.

Interpretation defines natural transformations between observational
functors.

Hence every admissible explanatory framework factors through the categorical
structure generated by autonomous systems and their admissibility-preserving
morphisms.

\end{proof}

This chapter completes a gradual shift in perspective that has developed
throughout the monograph.

We began with individual states.

We then introduced interfaces, repair, continuation, distinction geometry,
reachability, and local reconstruction.

Category theory reveals these constructions to be manifestations of a single
principle:

the mathematics of autonomous systems is fundamentally the mathematics of
structure-preserving transformations.

Objects exist only insofar as they participate in morphisms.

Identity exists only insofar as those morphisms preserve admissible
continuation.

The next chapter develops this observation into a unified theory of recursive
scientific discovery, showing that mathematical research itself is an
admissibility-preserving trajectory through a category of evolving theories.

\chapter{Scientific Discovery as Recursive Navigation}

\section{Discovery Is Not Creation}

Scientific discourse frequently describes discovery as though new knowledge
appeared suddenly through inspiration, intuition, or creativity.

Such descriptions correctly characterize the subjective experience of
discovery while obscuring its mathematical structure.

The present framework adopts a different perspective.

Scientific discovery is not the creation of new reality.

Neither is it the arbitrary invention of conceptual systems.

Instead, discovery consists of identifying admissible trajectories through
the space of possible theories.

The scientist does not create truth.

The scientist discovers navigable paths within an existing constraint
geometry.

\section{Theory Space}

Let

\[
\mathfrak{T}
\]

denote the collection of scientific theories under consideration.

Each element

\[
T\in\mathfrak{T}
\]

is itself an autonomous representational system consisting of

\begin{enumerate}
\item definitions,
\item propositions,
\item explanatory principles,
\item admissibility conditions,
\item observational commitments.
\end{enumerate}

Consequently,

\[
\mathfrak{T}
\]

is itself a configuration manifold.

Scientific development therefore becomes a dynamical process occurring upon
theory space rather than physical space.

\section{Theory Morphisms}

Theories evolve.

Definitions become refined.

Hypotheses become generalized.

Proofs become simplified.

This motivates the following construction.

\begin{definition}[Theory Morphism]

A theory morphism

\[
F:
T_1
\longrightarrow
T_2
\]

is an admissibility-preserving transformation satisfying

\begin{enumerate}
\item logical consistency,
\item preservation of established observational constraints,
\item preservation of previously verified internal theorems,
\item reduction of representational defect.
\end{enumerate}

\end{definition}

Scientific progress therefore consists of morphisms rather than isolated
theories.

Theories should be compared by their admissible transformations rather than
their historical chronology.

\section{Research Trajectories}

A research program determines a continuous path through theory space.

\begin{definition}[Research Trajectory]

A research trajectory is a continuous curve

\[
\Gamma:
[0,\infty)
\rightarrow
\mathfrak{T},
\]

whose image consists entirely of admissible theories.

\end{definition}

The parameter need not represent chronological time.

It may instead represent accumulated understanding, mathematical refinement,
or explanatory capacity.

Scientific maturity therefore corresponds to increasing regularity along the
trajectory rather than increasing age.

\section{The Discovery Functional}

Not every admissible movement through theory space represents scientific
progress.

We therefore introduce an optimization principle.

\begin{definition}[Discovery Functional]

The discovery functional is

\[
\mathcal{D}
:
\mathfrak{T}
\rightarrow
\mathbb{R},
\]

defined by

\[
\mathcal{D}(T)
=
\alpha
E(T)
-
\beta
C(T)
-
\gamma
\Delta(T)
+
\delta
R(T),
\]

where

\begin{itemize}
\item $E(T)$ measures explanatory power,
\item $C(T)$ measures unnecessary complexity,
\item $\Delta(T)$ measures reconstruction defect,
\item $R(T)$ measures admissible reachability.
\end{itemize}

The constants

\[
\alpha,\beta,\gamma,\delta>0
\]

determine the relative importance of each contribution.

\end{definition}

Scientific refinement seeks trajectories satisfying

\[
\mathcal{D}
(\Gamma(t))
\nearrow.
\]

Discovery therefore becomes constrained optimization over theory space.

\section{Recursive Refinement}

Scientific development rarely proceeds by replacement.

Instead it proceeds through recursive correction.

Let

\[
\mathcal{R}_{\mathfrak T}
:
\mathfrak T
\rightarrow
\mathfrak T
\]

denote the theory repair operator.

Repeated refinement generates

\[
T,
\quad
\mathcal R(T),
\quad
\mathcal R^2(T),
\quad
\ldots
\]

Each iteration removes inconsistencies while preserving explanatory
structure whenever possible.

Consequently mature theories should exhibit decreasing reconstruction defect
without sacrificing observational adequacy.

\section{Scientific Fixed Points}

Some theories become remarkably stable.

\begin{definition}[Scientific Fixed Point]

A theory

\[
T^\ast
\]

is a scientific fixed point whenever

\[
\mathcal R(T^\ast)
=
T^\ast,
\]

and every sufficiently small admissible perturbation satisfies

\[
\mathcal R(T^\ast+\varepsilon)
=
T^\ast
+
o(\varepsilon).
\]

\end{definition}

Fixed points do not represent absolute truth.

Rather, they represent locally stable explanatory structures.

Future refinement remains possible, but only through sufficiently large
extensions of observational or conceptual reach.

\section{The Geometry of Paradigm Change}

Scientific revolutions correspond neither to arbitrary replacement nor to
pure accumulation.

Instead they alter the topology of theory space itself.

New observational interfaces enlarge admissible regions.

Previously disconnected research trajectories become connected.

Formerly impossible continuations become reachable.

Thus paradigm shifts should be understood as topological transitions within
the manifold of admissible theories rather than purely sociological events.

\section{The Recursive Discovery Principle}

The previous constructions culminate in the principal theorem of this
chapter.

\begin{theorem}[Recursive Discovery Principle]

Scientific knowledge evolves through recursively repaired trajectories in
theory space that simultaneously decrease reconstruction defect, preserve
verified invariants, and increase admissible explanatory reachability.

\end{theorem}

\begin{proof}

Theory repair decreases inconsistency while preserving established
structure.

Admissibility prevents logically invalid transitions.

Reachability restricts refinement to realizable theoretical extensions.

Consequently every admissible research trajectory consists of recursively
repaired theories whose explanatory organization improves monotonically with
respect to the discovery functional.

\end{proof}

Scientific progress therefore admits a geometric interpretation.

The objective of research is not the accumulation of isolated facts.

Neither is it unrestricted theoretical invention.

The objective is the recursive navigation of theory space toward increasingly
stable regions possessing lower reconstruction defect, greater explanatory
compression, stronger invariant preservation, and larger admissible domains.

Accordingly, science itself becomes an autonomous repair process whose
trajectory is governed by the same principles that organize cognition,
computation, biological adaptation, and distributed intelligence.

\chapter{Compression, Parsimony, and Structural Minimality}

\section{The Principle of Structural Economy}

Scientific theories should not be judged solely by predictive accuracy.

Among theories exhibiting equivalent explanatory capability, preference should
be given to those possessing smaller structural complexity.

This observation is commonly associated with principles of parsimony.

Within the present framework, however, parsimony is not merely heuristic.

It arises naturally from admissible representation.

\section{Representational Complexity}

\begin{definition}[Structural Complexity]

Let

\[
T
\]

be an admissible theory.

Its structural complexity is

\[
K(T),
\]

defined as the minimal descriptive complexity required to reconstruct every
admissible consequence of

\[
T.
\]

\end{definition}

The precise choice of complexity measure is intentionally left abstract.

Kolmogorov complexity, minimum description length, proof complexity, or
categorical presentations may all instantiate

\[
K.
\]

\section{Explanatory Compression}

Scientific understanding is fundamentally compressive.

A successful theory replaces numerous independent observations by a smaller
family of governing principles.

\begin{definition}[Compression Ratio]

The explanatory compression of a theory is

\[
\mathcal{C}(T)
=
\frac{I_{\mathrm{obs}}}
     {K(T)},
\]

where

\[
I_{\mathrm{obs}}
\]

measures the informational content of the observations explained by the
theory.

\end{definition}

Larger values correspond to greater explanatory efficiency.

\section{Minimal Presentations}

Different formulations may describe identical mathematical structure.

Accordingly, minimality concerns equivalence classes rather than individual
expressions.

\begin{definition}[Minimal Presentation]

A presentation

\[
T^\star
\]

is minimal whenever

\[
K(T^\star)
=
\inf
\{
K(S)
:
S
\cong
T
\},
\]

where

\[
\cong
\]

denotes structural equivalence.

\end{definition}

Scientific elegance therefore reflects structural minimality rather than
stylistic preference.

\begin{theorem}[Structural Economy Principle]

Among observationally equivalent theories, maximal explanatory compression
corresponds to minimal structural complexity.

\end{theorem}

\chapter{Symmetry, Redundancy, and Canonical Form}

\section{Redundant Descriptions}

Distinct representations frequently encode identical mathematical objects.

Coordinates, notation, implementation, and language may differ while the
underlying structure remains unchanged.

Scientific progress therefore requires identifying canonical
representations.

\begin{definition}[Redundancy]

Two representations

\[
R_1,
R_2
\]

are redundant whenever they determine the same admissible equivalence class.

\end{definition}

Redundancy increases descriptive complexity without increasing explanatory
power.

\section{Canonical Forms}

A canonical representation is one that minimizes redundancy while preserving
every invariant relevant to admissibility.

Canonicalization is therefore an admissibility-preserving projection onto a
minimal representative of an equivalence class.

\begin{proposition}

Every reduction in representational redundancy strictly decreases
reconstruction ambiguity while preserving invariant structure.

\end{proposition}

\chapter{The Principle of Recursive Scientific Realism}

\section{Reality Through Repair}

Scientific realism is frequently interpreted as the claim that successful
theories correspond to external reality.

The present framework adopts a weaker and more constructive position.

A theory is accepted not because it has been proven final, but because it
admits indefinitely continued repair while preserving an expanding family of
invariants.

\begin{definition}[Recursively Realizable Theory]

A theory

\[
T
\]

is recursively realizable whenever there exists an infinite admissible
sequence

\[
T_0,
T_1,
T_2,
\ldots
\]

such that

\[
T_0=T,
\]

each transition is an admissible repair, and every verified invariant of

\[
T_n
\]

remains preserved in

\[
T_{n+1}.
\]

\end{definition}

\begin{theorem}[Recursive Scientific Realism]

Scientific progress is characterized neither by convergence toward immutable
certainty nor by unrestricted conceptual revision.

Rather, it consists of recursively realizable trajectories preserving an
expanding family of experimentally and mathematically verified invariants.

\end{theorem}

The objective of science is therefore not perfection.

It is indefinite admissible continuation.

\newpage

\begin{thebibliography}{99}

%%%%%%%%%%%%%%%%%%%%%%%%%%%%%%%%%%%%%%%%%%%%%%%%%%%%%%%%%%%%%%
%% Foundations
%%%%%%%%%%%%%%%%%%%%%%%%%%%%%%%%%%%%%%%%%%%%%%%%%%%%%%%%%%%%%%

\bibitem{MacLane}
Saunders Mac Lane,
\emph{Categories for the Working Mathematician},
2nd ed.,
Springer,
1998.

\bibitem{Awodey}
Steve Awodey,
\emph{Category Theory},
2nd ed.,
Oxford University Press,
2010.

\bibitem{LawvereSchanuel}
F. William Lawvere and Stephen Schanuel,
\emph{Conceptual Mathematics},
2nd ed.,
Cambridge University Press,
2009.

%%%%%%%%%%%%%%%%%%%%%%%%%%%%%%%%%%%%%%%%%%%%%%%%%%%%%%%%%%%%%%
%% Topology and Geometry
%%%%%%%%%%%%%%%%%%%%%%%%%%%%%%%%%%%%%%%%%%%%%%%%%%%%%%%%%%%%%%

\bibitem{Munkres}
James R. Munkres,
\emph{Topology},
2nd ed.,
Prentice Hall,
2000.

\bibitem{Lee}
John M. Lee,
\emph{Introduction to Topological Manifolds},
2nd ed.,
Springer,
2011.

\bibitem{LeeSmooth}
John M. Lee,
\emph{Introduction to Smooth Manifolds},
2nd ed.,
Springer,
2012.

\bibitem{GuilleminPollack}
Victor Guillemin and Alan Pollack,
\emph{Differential Topology},
AMS Chelsea,
2010.

%%%%%%%%%%%%%%%%%%%%%%%%%%%%%%%%%%%%%%%%%%%%%%%%%%%%%%%%%%%%%%
%% Differential Geometry
%%%%%%%%%%%%%%%%%%%%%%%%%%%%%%%%%%%%%%%%%%%%%%%%%%%%%%%%%%%%%%

\bibitem{doCarmo}
Manfredo P. do Carmo,
\emph{Differential Geometry of Curves and Surfaces},
Prentice Hall,
1976.

\bibitem{Spivak}
Michael Spivak,
\emph{A Comprehensive Introduction to Differential Geometry},
Publish or Perish,
1999.

%%%%%%%%%%%%%%%%%%%%%%%%%%%%%%%%%%%%%%%%%%%%%%%%%%%%%%%%%%%%%%
%% Dynamical Systems
%%%%%%%%%%%%%%%%%%%%%%%%%%%%%%%%%%%%%%%%%%%%%%%%%%%%%%%%%%%%%%

\bibitem{Hirsch}
Morris Hirsch,
Stephen Smale,
Robert Devaney,
\emph{Differential Equations, Dynamical Systems, and an Introduction to Chaos},
Academic Press,
2012.

\bibitem{Katok}
Anatole Katok and Boris Hasselblatt,
\emph{Introduction to the Modern Theory of Dynamical Systems},
Cambridge University Press,
1995.

%%%%%%%%%%%%%%%%%%%%%%%%%%%%%%%%%%%%%%%%%%%%%%%%%%%%%%%%%%%%%%
%% Optimization
%%%%%%%%%%%%%%%%%%%%%%%%%%%%%%%%%%%%%%%%%%%%%%%%%%%%%%%%%%%%%%

\bibitem{Boyd}
Stephen Boyd and Lieven Vandenberghe,
\emph{Convex Optimization},
Cambridge University Press,
2004.

\bibitem{Rockafellar}
R. Tyrrell Rockafellar,
\emph{Convex Analysis},
Princeton University Press,
1970.

%%%%%%%%%%%%%%%%%%%%%%%%%%%%%%%%%%%%%%%%%%%%%%%%%%%%%%%%%%%%%%
%% Information Theory
%%%%%%%%%%%%%%%%%%%%%%%%%%%%%%%%%%%%%%%%%%%%%%%%%%%%%%%%%%%%%%

\bibitem{Shannon}
Claude E. Shannon,
``A Mathematical Theory of Communication,''
\emph{Bell System Technical Journal},
27 (1948),
379--423, 623--656.

\bibitem{CoverThomas}
Thomas Cover and Joy Thomas,
\emph{Elements of Information Theory},
2nd ed.,
Wiley,
2006.

%%%%%%%%%%%%%%%%%%%%%%%%%%%%%%%%%%%%%%%%%%%%%%%%%%%%%%%%%%%%%%
%% Probability
%%%%%%%%%%%%%%%%%%%%%%%%%%%%%%%%%%%%%%%%%%%%%%%%%%%%%%%%%%%%%%

\bibitem{Billingsley}
Patrick Billingsley,
\emph{Probability and Measure},
3rd ed.,
Wiley,
1995.

%%%%%%%%%%%%%%%%%%%%%%%%%%%%%%%%%%%%%%%%%%%%%%%%%%%%%%%%%%%%%%
%% Graph Theory
%%%%%%%%%%%%%%%%%%%%%%%%%%%%%%%%%%%%%%%%%%%%%%%%%%%%%%%%%%%%%%

\bibitem{Diestel}
Reinhard Diestel,
\emph{Graph Theory},
5th ed.,
Springer,
2017.

%%%%%%%%%%%%%%%%%%%%%%%%%%%%%%%%%%%%%%%%%%%%%%%%%%%%%%%%%%%%%%
%% Algebra
%%%%%%%%%%%%%%%%%%%%%%%%%%%%%%%%%%%%%%%%%%%%%%%%%%%%%%%%%%%%%%

\bibitem{LangAlgebra}
Serge Lang,
\emph{Algebra},
Revised 3rd ed.,
Springer,
2002.

%%%%%%%%%%%%%%%%%%%%%%%%%%%%%%%%%%%%%%%%%%%%%%%%%%%%%%%%%%%%%%
%% Functional Analysis
%%%%%%%%%%%%%%%%%%%%%%%%%%%%%%%%%%%%%%%%%%%%%%%%%%%%%%%%%%%%%%

\bibitem{RudinFunctional}
Walter Rudin,
\emph{Functional Analysis},
2nd ed.,
McGraw--Hill,
1991.

%%%%%%%%%%%%%%%%%%%%%%%%%%%%%%%%%%%%%%%%%%%%%%%%%%%%%%%%%%%%%%
%% Logic
%%%%%%%%%%%%%%%%%%%%%%%%%%%%%%%%%%%%%%%%%%%%%%%%%%%%%%%%%%%%%%

\bibitem{Enderton}
Herbert Enderton,
\emph{A Mathematical Introduction to Logic},
2nd ed.,
Academic Press,
2001.

%%%%%%%%%%%%%%%%%%%%%%%%%%%%%%%%%%%%%%%%%%%%%%%%%%%%%%%%%%%%%%
%% Sheaf Theory
%%%%%%%%%%%%%%%%%%%%%%%%%%%%%%%%%%%%%%%%%%%%%%%%%%%%%%%%%%%%%%

\bibitem{Bredon}
Glen E. Bredon,
\emph{Sheaf Theory},
2nd ed.,
Springer,
1997.

\bibitem{Kashiwara}
Masaki Kashiwara and Pierre Schapira,
\emph{Categories and Sheaves},
Springer,
2006.

%%%%%%%%%%%%%%%%%%%%%%%%%%%%%%%%%%%%%%%%%%%%%%%%%%%%%%%%%%%%%%
%% Computation
%%%%%%%%%%%%%%%%%%%%%%%%%%%%%%%%%%%%%%%%%%%%%%%%%%%%%%%%%%%%%%

\bibitem{Sipser}
Michael Sipser,
\emph{Introduction to the Theory of Computation},
3rd ed.,
Cengage,
2013.

%%%%%%%%%%%%%%%%%%%%%%%%%%%%%%%%%%%%%%%%%%%%%%%%%%%%%%%%%%%%%%
%% Control Theory
%%%%%%%%%%%%%%%%%%%%%%%%%%%%%%%%%%%%%%%%%%%%%%%%%%%%%%%%%%%%%%

\bibitem{Bertsekas}
Dimitri Bertsekas,
\emph{Dynamic Programming and Optimal Control},
Athena Scientific,
2017.

%%%%%%%%%%%%%%%%%%%%%%%%%%%%%%%%%%%%%%%%%%%%%%%%%%%%%%%%%%%%%%
%% Scientific Method
%%%%%%%%%%%%%%%%%%%%%%%%%%%%%%%%%%%%%%%%%%%%%%%%%%%%%%%%%%%%%%

\bibitem{Popper}
Karl Popper,
\emph{The Logic of Scientific Discovery},
Routledge,
2002.

\bibitem{Lakatos}
Imre Lakatos,
\emph{The Methodology of Scientific Research Programmes},
Cambridge University Press,
1978.

\bibitem{Kuhn}
Thomas S. Kuhn,
\emph{The Structure of Scientific Revolutions},
4th ed.,
University of Chicago Press,
2012.

%%%%%%%%%%%%%%%%%%%%%%%%%%%%%%%%%%%%%%%%%%%%%%%%%%%%%%%%%%%%%%
%% Systems Theory
%%%%%%%%%%%%%%%%%%%%%%%%%%%%%%%%%%%%%%%%%%%%%%%%%%%%%%%%%%%%%%

\bibitem{vonBertalanffy}
Ludwig von Bertalanffy,
\emph{General System Theory},
George Braziller,
1968.

%%%%%%%%%%%%%%%%%%%%%%%%%%%%%%%%%%%%%%%%%%%%%%%%%%%%%%%%%%%%%%
%% Cybernetics
%%%%%%%%%%%%%%%%%%%%%%%%%%%%%%%%%%%%%%%%%%%%%%%%%%%%%%%%%%%%%%

\bibitem{Wiener}
Norbert Wiener,
\emph{Cybernetics},
MIT Press,
1961.


\end{thebibliography}

\end{document}
